% chapters/ch15_apcso_guide.tex
% Chapter 15: APCSO Operational Guide
% Source: NRMO_v5_updated.pdf, Chapter 15 (lines 8179-8424 of v5_layout.txt)
% Content overlaps with Section 12.2 (APCSO formal theory) and 12.3 (applied guide)

\chapter{APCSO Operational Guide --- Layered View}
\label{ch:apcso-guide}

\nrmobox{Document role}{
This chapter presents the layered APCSO architecture as an
operational reference. It is the third complementary view of APCSO,
following the formal theory (Section~\ref{sec:apcso-formal}) and the
applied guide (Section~\ref{sec:apcso-applied-guide}).
}

% =====================================================================
\section{APCSO Positioning}
\label{sec:apcso-positioning}

\paragraph{APCSO is not.}
APCSO is \textbf{not} an operation, OS, or replacement for any
existing layer. Its purpose is solely to optimise the
\emph{structure} of the choice set under irreversibility constraints.

\paragraph{APCSO is.}
APCSO is a layer that consults Passive Pattern (observation),
evaluates options, and outputs an unordered three-element proposal
set. It does not select; it does not execute; it does not rank.

% =====================================================================
\section{Layer Structure}
\label{sec:apcso-layer-structure}

\paragraph{Passive Pattern (observation layer).}
Observation; evaluation / classification. Shutdown via HALT or HOLD.

\paragraph{APCSO (judgement layer).}
Generates the choice set under conditions: exactly three options;
at least one hold or exit; ruin upper-bound $\le \varepsilon$ for
all.

\paragraph{Human / Governance (selection layer).}
$\varepsilon$ supply, selection, and accountability. The human
selects; the governance authority supplies the $\varepsilon$.

% =====================================================================
\section{Operating Conditions}
\label{sec:apcso-operating-conditions}

APCSO activates when:

\begin{itemize}
\item \textbf{Passive Pattern Gate clears} (state is SAFE).
\item \textbf{Irreversibility risk is non-trivial} (one or more
options approach the irreversible boundary).
\item \textbf{Non-zero ruin risk exists} (else the Zero-Risk Fast
Path applies; see Section~\ref{sec:zero-risk-fast-path}).
\item \textbf{Binary framing or forced-choice pressure detected}
(APCSO's purpose is to break this framing).
\end{itemize}

\noindent
APCSO does not produce action.

% =====================================================================
\section{Cross-References}

The complete APCSO theory and applied guide are in
Chapter~\ref{ch:operational-refinements}:

\begin{itemize}
\item Formal theory: Section~\ref{sec:apcso-formal}.
\item Applied guide and operational checklist:
Section~\ref{sec:apcso-applied-guide}.
\item Implementation mapping table:
Table~\ref{tab:apcso-implementation}.
\end{itemize}

% =====================================================================
\section*{Connections to Other Parts}

\begin{itemize}
\item APCSO sits between Passive Pattern (Section~\ref{sec:passive-pattern-paper})
and NRMO Core's admissibility evaluation
(Section~\ref{sec:foundations-3-1-definition}).
\item The proposal-set formal definition is Equation~\ref{eq:apcso-proposal-set}
(Section~\ref{sec:apcso-formal}).
\item Layer roles map to the governance--execution separation
invariant (Invariant~\ref{inv:gov-exec-sep}).
\end{itemize}
